%----------------------------------------------------------------------
\section{Discussion}
\label{sec:analysis}
%----------------------------------------------------------------------

\paragraph{The objective, not the architecture.}
The same backbone, the same corpus and the same number of steps produce either a forecaster that
improves with training or one that degrades with training, depending only on what the loss measures.
That is a stronger statement than "our loss is better": it says the usual practice of continuing a
vision model's native objective on rendered series is not a neutral default but an active liability
at scale, and that the community's instinct to buy accuracy with more pretraining will be punished
here. It also explains why the visual prior found by prior work looks small \citep{svtime}: when the
read-out is outside the loop, most of the capacity being transferred is being spent on things the
read-out discards.

\paragraph{What a video prior actually supplies.}
Two independent observations suggest the answer is long-range temporal context rather than anything
specific to video semantics. First, at 16-frame clips the video initialisation beats the image
initialisation by $3$--$4\%$, and at 32 frames the two converge (Appendix~\ref{app:16f}); lengthen
the window and part of the prior becomes redundant. Second, prior anchoring is essential at 16
frames and inert at 32. Both point at the same thing: the prior is worth most exactly where the
input does not already contain the structure. The matched-size comparison at a short budget
(Table~\ref{tab:backbone}) is where the video prior is visible, and the random-init control shows it
is real rather than an artefact of capacity.

We are careful about what this does not say. It does not say video backbones beat image backbones in
general; at 20k steps under the aligned objective the gap closes. It does not say predictive world
models beat reconstructive ones; at matched size they tie. The defensible claim is narrower and, we
think, more useful: a prior over dynamics transfers to forecasting, several pretraining objectives
produce one, and the objective used during transfer decides whether any of it survives.

\paragraph{Why the value-space term grows with horizon.}
At short horizons a correct completion is largely determined by the visible context, and pixel error
and value error rank candidate completions similarly. As the horizon grows, more of the prediction
is generated rather than interpolated, and the two metrics diverge: a prediction can be smooth,
plausible and pixel-cheap while being wrong by a constant offset in value. The ablation matches this
reading, with the gap growing from $5\%$ at $H{=}96$ to $8\%$ at $H{=}720$.

\paragraph{Limitations.}
Mid-to-long horizons on the hourly ETT sets remain behind VisionTS and VisionTS++, and the cause is
structural: our context is capped at 32 frames of at most four periods, so at $H{=}720$ with
$P{=}24$ we trade resolution for coverage. Mean MAE ties rather than beats the strongest baseline
($0.328$ against $0.326$), which is worth stating plainly since we lead with MSE. The evaluation is
the six-dataset LSF protocol only; GIFT-Eval \citep{gifteval} and Monash \citep{monash} results are
not yet in. Our budget comparison is in pretraining windows and observations, not FLOPs. The main
result uses one backbone family, with the V-JEPA arms serving as 5k-step controls rather than as
full models. Finally, the pretraining corpus contains Monash \texttt{traffic\_hourly}, so the
traffic benchmark cannot be added without rebuilding the corpus, and we exclude it rather than
report a contaminated number.
